\chapter{Wasserstein 距離 \texorpdfstring{$W_p$}{Wp}}
\label{ch:wasserstein-metrics}

本章を通して $(X,d)$ を Polish 空間とし，基準点 $x_0\in X$ を一つ固定する．

\section{定義}
\label{sec:wp-definition}

\begin{definition}{有限モーメントをもつ確率測度}{gs-p-space}
  $1\leq p<\infty$ に対し
  \[
    \Pp{p}(X)
    :=
    \left\{
      \mu\in\Prob(X)
      \;\middle|\;
      \int_X d(x,x_0)^p\,\d\mu(x)<\infty
    \right\}
  \]
  と定める．また
  \[
    \Pp{\infty}(X)
    :=
    \left\{\mu\in\Prob(X)\mid \spt\mu\text{ は有界}\right\}
  \]
  とおく．三角不等式より $\Pp{p}(X)$ は基準点 $x_0$ の選び方によらない．
  これは $\Wass_p$ を有限値の距離として扱うための定義域である．
\end{definition}

\begin{definition}{結合と Wasserstein 距離}{gs-wasserstein}
  $\mu,\nu\in\Prob(X)$ の\textbf{結合}の集合を
  \[
    \Couplings(\mu,\nu)
    :=
    \left\{
      \pi\in\Prob(X\times X)
      \;\middle|\;
      (\mathrm{pr}_1)\pushforward\pi=\mu,\ 
      (\mathrm{pr}_2)\pushforward\pi=\nu
    \right\}
  \]
  とする．$1\leq p<\infty$ と $\mu,\nu\in\Pp{p}(X)$ に対し
  \[
    \Wass_p(\mu,\nu)
    :=
    \inf_{\pi\in\Couplings(\mu,\nu)}
    \left(
      \int_{X\times X}d(x,y)^p\,\d\pi(x,y)
    \right)^{1/p},
  \]
  また $\mu,\nu\in\Pp{\infty}(X)$ に対し
  \[
    \Wass_\infty(\mu,\nu)
    :=
    \inf_{\pi\in\Couplings(\mu,\nu)}
    \esssup_{(x,y)\sim\pi}d(x,y)
  \]
  と定める．
\end{definition}

$p<\infty$ のとき，直積結合 $\mu\otimes\nu$ と
\[
  d(x,y)^p
  \leq
  2^{p-1}\bigl(d(x,x_0)^p+d(y,x_0)^p\bigr)
\]
を用いれば，$\mu,\nu\in\Pp{p}(X)$ に対して $\Wass_p(\mu,\nu)<\infty$ が分かる．

\section{最適結合の存在}
\label{sec:optimal-coupling}

\begin{lemma}{輸送コストの下半連続性}{gs-cost-lsc}
  $\pi_n\rightharpoonup\pi$ ならば，$1\leq p<\infty$ に対し
  \[
    \int d(x,y)^p\,\d\pi(x,y)
    \leq
    \liminf_{n\to\infty}
    \int d(x,y)^p\,\d\pi_n(x,y).
  \]
\end{lemma}

\begin{proof}
  $(x,y)\mapsto d(x,y)^p$ は非負かつ連続なので，
  Portmanteau の定理から従う．
\end{proof}

\begin{proposition}{最適結合の存在（原論文 Proposition 1）}{gs-optimal-coupling}
  $1\leq p\leq\infty$ とする．
  $\mu,\nu$ が $\Wass_p$ の定義域に属するならば，
  下限を達成する $\pi^\star\in\Couplings(\mu,\nu)$ が存在する．
\end{proposition}

\begin{proof}
  $p<\infty$ とする．最小化列
  $(\pi_n)\subset\Couplings(\mu,\nu)$ を取る．
  固定された周辺分布 $\mu,\nu$ は tight なので，
  $\Couplings(\mu,\nu)$ も tight である．
  Prokhorov の定理により，部分列を取れば
  $\pi_{n_k}\rightharpoonup\pi^\star\in\Couplings(\mu,\nu)$ とできる．
  補題~\ref{lem:gs-cost-lsc}より $\pi^\star$ は下限を達成する．

  $p=\infty$ では
  \[
    \esssup_{\pi_n}d
    \leq
    \Wass_\infty(\mu,\nu)+\frac1n
  \]
  を満たす結合を取り，弱収束部分列 $\pi_{n_k}\rightharpoonup\pi^\star$ を取る．
  各 $m$ に対し，十分大きな $k$ で
  $\pi_{n_k}(\{d\leq\Wass_\infty+1/m\})=1$ である．
  この集合は閉なので，Portmanteau の定理から
  $\pi^\star(\{d\leq\Wass_\infty+1/m\})=1$．
  $m\to\infty$ とすれば $\pi^\star$ は下限を達成する．
\end{proof}

\section{Proposition 2：距離性}
\label{sec:metric-property}

\begin{lemma}{結合の貼り合わせ}{gs-gluing}
  $\pi_{12}\in\Couplings(\mu_1,\mu_2)$ と
  $\pi_{23}\in\Couplings(\mu_2,\mu_3)$ に対し，
  $X^3$ 上の確率測度 $\boldsymbol{\pi}$ で
  $(x_1,x_2)$ 周辺分布が $\pi_{12}$，
  $(x_2,x_3)$ 周辺分布が $\pi_{23}$ となるものが存在する．
\end{lemma}

\begin{proof}
  正則条件付き確率を用いて
  \[
    \pi_{12}(\d x_1,\d x_2)
    =K_{1|2}(x_2,\d x_1)\mu_2(\d x_2),
    \quad
    \pi_{23}(\d x_2,\d x_3)
    =\mu_2(\d x_2)K_{3|2}(x_2,\d x_3)
  \]
  と分解し，
  \[
    \boldsymbol{\pi}(\d x_1,\d x_2,\d x_3)
    :=
    K_{1|2}(x_2,\d x_1)\mu_2(\d x_2)K_{3|2}(x_2,\d x_3)
  \]
  とおけばよい．
\end{proof}

\begin{proposition}{Wasserstein 距離は距離である（原論文 Proposition 2）}{gs-wp-metric}
  $1\leq p\leq\infty$ に対し，
  $\Wass_p$ は $\Pp{p}(X)$ 上の距離である．
\end{proposition}

\begin{proof}
  非負性と対称性は定義から明らかである．
  $\Wass_p(\mu,\nu)=0$ なら，命題~\ref{prop:gs-optimal-coupling}の
  最適結合は対角集合に集中するので $\mu=\nu$ である．

  三角不等式を示す．$\mu_1,\mu_2,\mu_3\in\Pp{p}(X)$ とし，
  $\pi_{12},\pi_{23}$ を最適結合とする．
  補題~\ref{lem:gs-gluing}でこれらを貼り合わせる．
  $p<\infty$ なら，距離 $d$ の三角不等式と Minkowski の不等式から
  \[
    \begin{aligned}
      \Wass_p(\mu_1,\mu_3)
      &\leq \norm{d(x_1,x_3)}_{L^p(\boldsymbol{\pi})} \\
      &\leq \norm{d(x_1,x_2)+d(x_2,x_3)}_{L^p(\boldsymbol{\pi})} \\
      &\leq \Wass_p(\mu_1,\mu_2)+\Wass_p(\mu_2,\mu_3).
    \end{aligned}
  \]
  $p=\infty$ でも同じ議論を本質的上限に適用すればよい．
\end{proof}
